\documentclass[10pt,a4paper]{amsart}
\usepackage[margin=19mm]{geometry}
\usepackage{amsmath,amssymb,mathtools}
\usepackage[scr=rsfs,cal=euler]{mathalfa}
\usepackage{xfrac}
\usepackage[unicode,hidelinks,bookmarks=false]{hyperref}
\newcommand{\ie}{i.e.\ }
\newcommand{\cf}{cf.\ }
\newcommand{\resp}{respectively\ }
\newcommand{\Subsection}{\subsection}
\providecommand{\nobreakdash}{\nobreak\hbox{-}}
\setlength{\emergencystretch}{2em}
\multlinegap0pt
\usepackage{tikz}
\usepackage{amssymb}
\newtheorem{lemma}{Lemma}
\title{Directed homological and cohomological operations}
\author{Eric Goubault}
\date{Source: arXiv:2603.11967v1 (CC BY 4.0)}
\begin{document}
\maketitle

\noindent\emph{Source and licence.} Directed homological and cohomological operations. Version arXiv:2603.11967v1 (CC BY 4.0), distributed by arXiv under the Creative Commons Attribution 4.0 International licence, http://creativecommons.org/licenses/by/4.0/, as recorded in the arXiv licence metadata for this exact version and shown on the abstract page https://arxiv.org/abs/2603.11967v1. Full licence notice: \texttt{licenses/arxiv-2603.11967-CC-BY-4.0.txt}.\par\smallskip

\noindent\textit{Editorial scope.} Counted unit: the article's lemma lem:diffalg (source0.tex lines 554--560). The article's complete original proof of it is delivered with it (source0.tex lines 562--591, one \texttt{proof} environment of 30 lines). No mathematics is written, repaired or re-derived: the statement and the proof are the article's own lines, in the article's own notation, and no retained line is substituted or re-flowed.  No further source-local context is needed: every cross-reference of every retained line resolves inside the two delivered spans.  Nothing was omitted from any retained span, so the package carries no omission record.  No retained line contains a citation, so the wrapper carries no bibliography and no bibliography entry.  No retained line contains a figure environment, an image-inclusion command or drawing code, so no image is delivered and the package declares no asset dependency.  Editorial formatting only, in addition: an AMS \texttt{amsart} preamble that restates the article's own declarations and macros used by the retained lines, namely the environments \texttt{lemma} on separate counters, together with the standard packages those lines need.  The retained environments print in this document as Lemma 1 in order of appearance, whereas the article prints the same items with the numbering of its own full document.  The complete native source of the article as distributed in the arXiv e-print for this exact version is bundled unchanged as \texttt{sources/196248/source0.tex}.
\begin{lemma}
\label{lem:diffalg}
    The boundary operator $\partial$ has the property that, for any $c$ $(i+1)$-cube path and $d$ any $(j+1)$-cube path:
$$
\partial(c\otimes d) = \partial(c)\otimes d +(-1)^{|c|} c\otimes \partial(d)
$$
\end{lemma}

\begin{proof}
Suppose first that $c$ and $d$ are cube chains

When $c$ and $d$ do not have matching ends, the equality above is trivial (everything is 0).

Suppose now that $c$ and $d$ have matching ends. 
Let $c=(c_1,\ldots,c_n)$ (so that $|c|=n$ is the length of $c$), and $d=(d_1,\ldots,d_m)$. Therefore we have $c\otimes d=(c_1,\ldots,c_n,d_1,\ldots,d_m)$.

Now: 

$$
\begin{array}{rcl}
{\scriptstyle \partial (c\otimes d)} & {\scriptstyle =} &\sum\limits_{k=1}^{n+m} \sum\limits_{r=1}^{n_k-1} \sum\limits_{\scriptscriptstyle \mathcal{I} \subseteq\{0,\ldots,n_k-1\}: \ \mid \mathcal{I} \mid=r}
{\scriptstyle (-1)^{\scriptscriptstyle n_1+\ldots+n_{k-1}+k+r+1} sgn(\mbox{$\scriptstyle \mathcal{I}$}) d_{k,\mbox{$\scriptscriptstyle \mathcal{I}$}}(c\otimes d)} \\
& {\scriptstyle =} & \sum\limits_{k=1}^{n} \sum\limits_{r=1}^{n_k-1} \sum\limits_{\scriptscriptstyle \mathcal{I} \subseteq\{0,\ldots,n_k-1\}: \ \mid \mathcal{I} \mid=r}
{\scriptstyle (-1)^{\scriptscriptstyle n_1+\ldots+n_{k-1}+k+r+1} sgn(\mbox{$\scriptstyle \mathcal{I}$}) d_{k,\mbox{$\scriptscriptstyle \mathcal{I}$}}(c\otimes d)} \\
& & +\sum\limits_{k=n+1}^{m} \sum\limits_{r=1}^{n_k-1} \sum\limits_{\scriptscriptstyle \mathcal{I} \subseteq\{0,\ldots,n_k-1\}: \ \mid \mathcal{I} \mid=r} 
{\scriptstyle (-1)^{\scriptscriptstyle n_1+\ldots+n_{k-1}+k+r+1}
sgn(\mbox{$\scriptstyle \mathcal{I}$}) d_{k,\mbox{$\scriptscriptstyle \mathcal{I}$}}
 (c\otimes d)} \\ 
& {\scriptstyle =} & \sum\limits_{k=1}^{n} \sum\limits_{r=1}^{n_k-1} \sum\limits_{\scriptscriptstyle \mathcal{I} \subseteq\{0,\ldots,n_k-1\}: \ \mid \mathcal{I} \mid=r}
{\scriptstyle (-1)^{\scriptscriptstyle n_1+\ldots+n_{k-1}+k+r+1} sgn(\mbox{$\scriptstyle \mathcal{I}$}) d_{k,\mbox{$\scriptscriptstyle \mathcal{I}$}}(c)\otimes d} \\
& & +(-1)^n \sum\limits_{k=1}^{m} \sum\limits_{r=1}^{n_k-1} \sum\limits_{\scriptscriptstyle \mathcal{I} \subseteq\{0,\ldots,n_k-1\}: \ \mid \mathcal{I} \mid=r}
{\scriptstyle (-1)^{\scriptscriptstyle n_1+\ldots+n_{k-1}+k+r+1} sgn(\mbox{$\scriptstyle \mathcal{I}$}) c\otimes d_{k,\mbox{$\scriptscriptstyle \mathcal{I}$}}(d)} \\
& {\scriptstyle =} & {\scriptstyle \partial(c)\otimes d +(-1)^{|c|}c\otimes d}
\end{array}
$$
%
%The case when $c$ and $d$ are general elements of $R_i[X]$ and of $R_j[X]$ respectively stems out from bilinearity of the tensor operation. 
\end{proof}

\end{document}
